\section{Removing a zero tail and completing the proof}
\label{sec:zero-tail-completion}

Theorem~\ref{thm:integral-schedule} applies when the moving weight has
at most one zero coordinate.  We now supply the real existence input
needed when all three labels have a longer zero tail. The input is a
tensor occurrence for labels supported on a common initial coordinate
segment. We construct an occurrence for twice the truncated labels in
a smaller even spin group. Passing to real moments then removes the
factor of two. The integral schedule will be constructed for the
unscaled truncated labels, using the restriction of the original
column shape.

Write the coordinates of the labels as
\(\lambda=(L_1,\ldots,L_n)\), \(\mu=(M_1,\ldots,M_n)\), and
\(\xi=(X_1,\ldots,X_n)\). The reduction we shall prove is the
following.

\begin{proposition}[Removal of a zero tail]\label{prop:zero-tail}
Suppose \(\lambda,\mu,\xi\) are dominant integral \(D_n\)-weights,
all supported in coordinates \(1,\ldots,r\), where
\(1\le r\le n-2\).  If \(V(\xi)\) occurs in
\(V(\lambda)\otimes V(\mu)\), then
\[
 V(2\widehat\xi)\ \text{occurs in}\
 V(2\widehat\lambda)\otimes V(2\widehat\mu)
 \quad\text{for }D_{r+1},
\]
where
\(\widehat\lambda=(L_1,\ldots,L_r,0)\), and similarly for the
other two labels.
\end{proposition}

\subsection{A preprojective-module criterion}

Lusztig's semicanonical basis relates enveloping algebras to
preprojective representation varieties \cite{Lusztig00}. We use its
generic-component coordinates, together with their left and right
filtrations \cite{GLS05,FLX}, to turn tensor occurrence into
simultaneous bounds on the top and socle of one module. Those bounds
will let us change the two orthogonal leaves while controlling both
tensor factors.

Let \(\mathcal D\) be a finite simply laced Dynkin diagram, possibly
disconnected.  Choose an orientation of its edges and adjoin a reverse
arrow to every arrow.  A representation of the resulting doubled
quiver is \emph{preprojective} if, at each vertex, the signed sum of
the two-arrow return maps is zero.  Our sign convention assigns
\(uv\) a plus sign at \(j\) and \(vu\) a minus sign at \(i\), for an
oriented arrow \(u:i\longrightarrow j\) with reverse \(v\).
Write \(\Lambda(d)\) for the variety of these representations on fixed
spaces of dimension vector \(d=(d_i)\).

For \(M\in\Lambda(d)\), set
\[
 \operatorname{top}_i M
 =\dim\left(M_i\Big/\sum_{a:\,j\to i}\im M_a\right),
 \qquad
 \operatorname{soc}_i M
 =\dim\bigcap_{a:\,i\to j}\ker M_a .
\]
Both sums and intersections run over arrows of the doubled quiver.

\begin{proposition}[Tensor criterion]\label{prop:preprojective}
Let \(\lambda,\mu,\xi\) be dominant integral weights for \(\mathcal D\),
and suppose
\(\lambda+\mu-\xi=\sum_i d_i\alpha_i\) with \(d_i\in\Z_{\ge0}\).
Put \(\lambda_i=\lambda\cdot\alpha_i\) and
\(\mu_i=\mu\cdot\alpha_i\).  Then
\begin{equation}\label{tail:tensor-criterion}
 \begin{aligned}
 &\Hom\bigl(V(\lambda)\otimes V(\mu),V(\xi)\bigr)\ne0\\
 &\quad\Longleftrightarrow\quad
 \exists M\in\Lambda(d):
 \quad \operatorname{top}_i M\le\lambda_i,\quad
       \operatorname{soc}_i M\le\mu_i\quad\text{for every }i .
 \end{aligned}
\end{equation}
\end{proposition}

\begin{proof}
We recall precisely the semicanonical facts used here.  On the
preprojective varieties, convolution of invariant constructible
functions is defined by integration with compactly supported Euler
characteristic over submodules, with the submodule as the right
factor.  The constant function on the unique representation of
dimension \(e_i\) represents \(f_i\).  This identifies
\(U=U(\mathfrak n^-)\) with the algebra generated by these functions;
we identify the positive enveloping algebra in the cited conventions
with the negative one by the Chevalley isomorphism.
In each degree, generic evaluation on irreducible components is a
vector-space isomorphism
\begin{equation}\label{tail:generic-evaluation}
 U_d \longrightarrow
 \C^{\operatorname{Irr}\Lambda(d)} .
\end{equation}
Moreover, the image of \((f_i^pU)_d\) is the coordinate subspace
supported on components whose generic \(i\)-top is at least \(p\).
The image of \((Uf_i^p)_d\) is the analogous subspace for generic
\(i\)-socle.  Finite Dynkin type needs no additional nilpotence
condition.  For the constructible-function realization and generic
component coordinates, see
\cite[Proposition~3.1 and Sections~5.2--5.6]{GLS05}; for their
compatibility with the top and socle filtrations, see
\cite[Lemmas~5.1, 5.4, 5.6, and 5.9]{FLX}.

Here is the geometry behind the filtration statement, which also
fixes the left/right convention.  The preprojective equations
describe the union of conormals to the finitely many ordinary-quiver
orbits.  Its components have dimension
\[
 D(d)=\sum_{\{i,j\}\in\mathcal D}d_i d_j
\]
and are counted by the Kostant partitions of \(d\), hence by
\(\dim U_d\).  On the stratum with \(i\)-top exactly \(p\), remove the
top by replacing \(M_i\) with its incoming image, of dimension
\(a=d_i-p\).  The lower representation has \(i\)-top zero.
Conversely, put \(S=\sum_{j\sim i}d_j\).  After choosing the retained
\(a\)-plane, extensions add outgoing maps on the remaining
\(p\) dimensions with values in the kernel of the signed incoming
map, of dimension \(S-a\).  Their parameter count, including the
Grassmannian, is
\[
 p(S-a)+pa=pS=D(d)-D(d-pe_i).
\]
The other relations are unchanged, because a path entering vertex
\(i\) enters its retained image.  This bundle construction gives the
component correspondence for exact top \(p\).

At exact top \(p\), the submodule with quotient \(S_i^p\) is unique,
so multiplication by \(f_i^{(p)}=f_i^p/p!\) transfers the generic
value from the lower component.  It vanishes when the top is smaller
than \(p\).  Descending induction on the top removes any contributions
on larger-top components.  Together with induction on dimension and
the component count, this explains both
\eqref{tail:generic-evaluation} and its left filtration.
Transposition of the arrows gives the right filtration and socle.
Thus the same generic-coordinate basis is compatible with all the
subspaces just described.

By Lemma~\ref{lem:coefficient} and the highest-generator presentation
of \(V(\mu)\), the multiplicity in question is the dimension of
\begin{equation}\label{tail:double-quotient}
 U_d\Big/
 \left(
   \sum_i f_i^{\lambda_i+1}U+
   \sum_i U f_i^{\mu_i+1}
 \right)_d .
\end{equation}
Under \eqref{tail:generic-evaluation}, this quotient is nonzero
exactly when there is a component on which all the generic top and
socle bounds in \eqref{tail:tensor-criterion} hold.  Those bounds are
open rank conditions.  Their intersection on such a component is
nonempty; conversely, any point satisfying them lies on a component
where they hold generically.  This proves the criterion.
\end{proof}

\subsection{Bounds along the chain}

To shorten the diagram, we will retain an initial chain and replace
its continuation by two leaves. We first bound
the incoming-image annihilator and outgoing kernel at each possible
cut, on every level of the nilpotent filtration.

For \(D_n\), use simple roots
\(\alpha_i=e_i-e_{i+1}\) for \(i<n\), and
\(\alpha_n=e_{n-1}+e_n\).
Suppose that the occurrence in
Proposition~\ref{prop:preprojective} holds, and choose a module from
that criterion. Write \(p=n-1\).
Combine the two leaf spaces into
\[
 V_p=W=W_+\oplus W_-,
\]
following the trunk \(V_1,\ldots,V_{p-1}\), and put \(V_0=0\).
With arrows oriented to the right, denote the paired maps on edge
\(i\) by
\[
 u_i:V_{i-1}\longrightarrow V_i,\qquad
 v_i:V_i\longrightarrow V_{i-1}.
\]
Here \(u_p\) combines the two leaf maps and \(v_p\) combines their
reverses.  The relations are
\begin{equation}\label{tail:chain-relations}
 t_i=u_iv_i=v_{i+1}u_{i+1}\quad(i<p),\qquad
 t_p=u_pv_p,\qquad
 (t_p)_{++}=(t_p)_{--}=0 .
\end{equation}
Since \(t_1=0\), all \(t_i\) are nilpotent.  This follows inductively
from the fact that \(uv\) is nilpotent if and only if \(vu\) is
nilpotent.  With \(\Delta_j=L_j+M_j-X_j\), the dimensions are
\begin{equation}\label{tail:chain-dimensions}
 \dim V_i=\sum_{j\le i}\Delta_j\qquad(1\le i\le p).
\end{equation}
The last equality gives the sum of the two original leaf dimensions.
Top bounds on the trunk are \(L_i-L_{i+1}\), and the two leaf bounds
are \(L_p-L_n,L_p+L_n\); socle bounds use \(M\) in place of \(L\).

For a single edge, the nilpotent operator
\[
 \begin{pmatrix}0&v_i\\u_i&0\end{pmatrix}
 \quad\text{on }V_{i-1}\oplus V_i
\]
has a Jordan basis whose rows alternate between the two spaces.
Such a basis is obtained by choosing homogeneous top generators of
each Jordan length.  In one row, let \(h_L,h_R\) be the numbers of
vectors on the left and right, and let \(T,S\in\{0,1\}\) record
whether its top and bottom, respectively, lie on the right.  Then
\begin{equation}\label{tail:row-identity}
 h_R=h_L-1+T+S.
\end{equation}
For \(b\ge1\), define
\[
 q_i(b)=\dim\bigl(\ker t_i\cap\im t_i^{b-1}\bigr),
 \qquad q_0(b)=0,
\]
the number of \(t_i\)-rows of length at least \(b\).  Let
\(K_i^T(b)\) be the sum of \(T\) over edge-\(i\) rows with
\(h_R\ge b\), and let \(K_i^{\prime T}(b)\) use \(h_L\ge b\).
Define \(K_i^S,K_i^{\prime S}\) similarly.  These quantities do not
depend on a common choice of bases for adjacent edges.

\begin{lemma}[Threshold bounds]\label{tail:threshold-bounds}
For every \(1\le r\le p\) and \(b\ge1\),
\begin{equation}\label{tail:threshold-result}
 K_r^T(b)\le L_r+\frac{q_r(b)}2,\qquad
 K_r^S(b)\le M_r+\frac{q_r(b)}2 .
\end{equation}
\end{lemma}

\begin{proof}
Fix \(b\) and suppress it from the notation. For each vertex, write
\[
 F_i^*(b)=\ker t_i^*\cap\im(t_i^*)^{b-1}
\]
for its dual high-socle space. At a trunk vertex \(i<p\), the two
incoming maps give
\[
 \begin{aligned}
 \dim\bigl(\operatorname{Ann}(\im u_i)\cap F_i^*(b)\bigr)
   &=K_i^T,\\
 \dim\bigl(\operatorname{Ann}(\im v_{i+1})\cap F_i^*(b)\bigr)
   &=q_i-K_{i+1}^{\prime T}.
 \end{aligned}
\]
These identities are read directly from the alternating rows:
the left image fails to contain the right-hand top exactly when
\(T=1\), whereas the right image fails to contain a left-hand top
exactly when \(T=0\).
The intersection of the two annihilators is bounded by the full
vertex top.  The dimension inequality therefore gives
\begin{equation}\label{tail:trunk-inequality}
 K_i^T\le L_i-L_{i+1}+K_{i+1}^{\prime T}\qquad(i<p).
\end{equation}

At \(W\), the space \(F_p^*(b)\) splits into its two colors, since
\(t_p\) is off diagonal.  Let their dimensions be \(q_+,q_-\),
so \(q_++q_-=q_p\).  The left incoming-image annihilator has
intersection dimension \(K_p^T\) with \(F_p^*(b)\).  Its intersection
with either color is bounded by the corresponding leaf top, whence
\[
 K_p^T-q_-\le L_p+L_n,\qquad
 K_p^T-q_+\le L_p-L_n
\]
after naming the colors accordingly.  Adding proves
\begin{equation}\label{tail:leaf-inequality}
 K_p^T-\frac{q_p}{2}\le L_p.
\end{equation}
No equality of the old color dimensions is required.

For each edge, \eqref{tail:row-identity} also gives
\begin{equation}\label{tail:edge-comparison}
 K_i^{\prime T}-\frac{q_{i-1}}2
 \le K_i^T-\frac{q_i}2 .
\end{equation}
Indeed, the contribution of one row to right minus left is
\[
 \left(\mathbf1_{\{h_R\ge b\}}-\mathbf1_{\{h_L\ge b\}}\right)
 \left(T-\frac12\right).
\]
It is zero if the indicators agree.  If a row is gained on the
right, then \(T=S=1\); if it is lost, then \(T=S=0\).
In either case its contribution is \(1/2\).

Combining \eqref{tail:trunk-inequality} with
\eqref{tail:edge-comparison} for edge \(i+1\) gives
\[
 K_i^T-\frac{q_i}2
 \le L_i-L_{i+1}+K_{i+1}^T-\frac{q_{i+1}}2 .
\]
Starting from \eqref{tail:leaf-inequality} and descending proves the
top assertion. At a trunk vertex, on the primal high-socle space
\(\ker t_i\cap\im t_i^{b-1}\), the outgoing kernels
\(\ker v_i\) and \(\ker u_{i+1}\) have intersection dimensions
\(K_i^S\) and \(q_i-K_{i+1}^{\prime S}\), respectively.
Their common intersection is bounded by the vertex socle. The same
leaf estimate and row comparison, with \(T,L\) replaced by \(S,M\),
therefore prove the socle assertion.
\end{proof}

\subsection{Positioning subspaces in a nilpotent filtration}

At each possible cut, the threshold estimates now control the
incoming annihilator and outgoing kernel on every nilpotent
filtration level. We next determine how such profiles bound
intersections with a moving subspace, and which motions can be
induced by automorphisms commuting with the nilpotent operator.

\begin{lemma}[Generic filtered intersection]\label{lem:filtered-intersection}
Let
\[
 H=F_1\supseteq F_2\supseteq\cdots\supseteq F_{s+1}=0
\]
be a filtration of a finite-dimensional complex vector space, and
let \(P\subseteq\GL(H)\) be its stabilizer.  For subspaces \(A,C\)
write
\(a_b=\dim(A\cap F_b)\), \(c_b=\dim(C\cap F_b)\), and
\(f_b=\dim F_b\).  For \(g\) in a nonempty Zariski-open subset of \(P\),
\begin{equation}\label{tail:intersection-formula}
 \dim(A\cap gC)=
 \max\left(0,\max_{1\le b\le s}(a_b+c_b-f_b)\right).
\end{equation}
\end{lemma}

\begin{proof}
Each term on the right is a lower bound by the dimension inequality
inside \(F_b\).  To attain all the bounds, choose coordinates adapted
to the filtration.  A basis for a subspace with profile \(a_b\)
consists of \(a_b-a_{b+1}\) generic vectors with allowed coordinates
in \(F_b\), for each \(b\); the analogous description applies to
\(C\).  Such bases have the required profiles on a nonempty open set.

The coordinate supports of the combined list of basis vectors are
nested. This is the nested-support case of Hall's distinct-representatives
criterion \cite[Theorem~1]{Hall1935}. A maximum matching of these vectors to distinct allowed
coordinates leaves exactly
\[
 \max\left(0,\max_b(a_b+c_b-f_b)\right)
\]
unmatched vectors.  Indeed any collection of vectors has, as its
union of allowed coordinates, its largest allowed filtration level.
For level \(F_b\), at most \(a_b+c_b\) vectors are constrained to that
level or deeper.  These are all the matching constraints, and they
are attained by matching from the deepest level outward.  A matching
gives a nonzero minor, so the generic rank of the combined list has
exactly this deficiency.  Since each separate list is independent,
the deficiency is its intersection dimension.

Finally, \(P\) is transitive on subspaces with a prescribed filtration
profile: extend a filtration-adapted basis of the subspace to one of
\(H\).  We may therefore hold \(A\) fixed and vary only \(C\), which
gives \eqref{tail:intersection-formula}.  Maximal rank is an open
condition.
\end{proof}

We next recall the centralizer's action on the socle filtration.
This is dual to the quotient-filtration statement in
\cite[Section~3.6]{Fresse09}; we include the short Jordan-basis proof.

\begin{lemma}[Centralizer on the socle]\label{tail:centralizer}
Let \(t\) be nilpotent on a complex vector space \(V\), and put
\[
 F_b(t)=\ker t\cap\im t^{b-1}\qquad(b\ge1).
\]
The group
\[
 C(t)=\{g\in\GL(V):gt=tg\}
\]
induces the full stabilizer of this filtration on \(\ker t\).
It also induces the full stabilizer of the corresponding filtration
on \(\ker t^*\subseteq V^*\), under the contragredient action.
The algebraic group \(C(t)\) is irreducible.
\end{lemma}

\begin{proof}
Write \(V\) as a sum of Jordan row modules
\(\C[t]/(t^a)\).  A module map from a row of length \(a\) to a row of
length \(b\) can send its bottom nontrivially to the target bottom
exactly when \(a\le b\).  In that case, sending its top to a scalar
multiple of \(t^{b-a}\) times the target top gives any desired bottom
coefficient.  These are precisely the matrix entries permitted by
the filtration \(F_b(t)\).  Every filtration-preserving bottom matrix
therefore lifts to a commuting endomorphism.  If its action on
\(\ker t\) is invertible, the lift is invertible: a nonzero
\(t\)-stable kernel would contain a nonzero vector in \(\ker t\).
This proves the first assertion.  Apply the same argument to \(t^*\)
and use inverse transpose for the dual assertion.

Finally, \(C(t)\) is the determinant-nonzero open subset of the vector
space \(\End_{\C[t]}(V)\), so it is irreducible.
\end{proof}

\subsection{The doubling construction}

The threshold inequalities control every filtration level at which
a new leaf bound can fail. After doubling they read
\(2K_r^T(b)-q_r(b)\le2L_r\) and
\(2K_r^S(b)-q_r(b)\le2M_r\).
Lemma~\ref{lem:filtered-intersection} will convert these inequalities
into top and socle bounds at both new leaves, once their spaces are
chosen with equal filtration profiles.

The doubled occurrence in Proposition~\ref{prop:zero-tail} is also
accessible through existing folding and stability results.
Hong--Shen's comparison for the fork automorphism carries occurrence
for invariant $D_n$ labels to saturated-cone membership for the
corresponding $C_{n-1}$ labels
\cite[Proposition~1.4(3)]{HS15}. Zero-padding stability of the
saturated cone in type $C$
\cite[Theorem~1.1]{GORY22}, followed by Sam's factor-two theorem
\cite[Theorem~1.1]{Sam12} and the reverse occurrence comparison
\cite[Proposition~1.4(1)]{HS15}, gives the doubled occurrence.
At the bottom ranks one uses $D_3=A_3$ and
$D_2=A_1\times A_1$. We instead give the explicit module construction,
which realizes the reduction while retaining the two sets of bounds
at the new leaves.

\begin{proof}[Proof of Proposition~\ref{prop:zero-tail}]
Occurrence implies that the simple-root coefficients of
\(\lambda+\mu-\xi\) are nonnegative integers, so choose a module
from Proposition~\ref{prop:preprojective}.
Keep the chain through \(V_r\) and take two copies of it.
The doubled space
\(\widetilde W=V_r\oplus V_r\) will be the new combined leaves.
Write \(\widetilde t=t_r\oplus t_r\),
\(\widetilde u=u_r\oplus u_r\), and
\(\widetilde v=v_r\oplus v_r\).

Color every Jordan row of \(t_r\) alternately, and give its duplicate
the opposite colors.  Thus \(\widetilde t\) is off diagonal for the
new splitting \(\widetilde W=\widetilde W_+\oplus\widetilde W_-\).
Each color has dimension \(\dim V_r\). A row of length \(a\) and
its duplicate contribute one bottom vector of each color to every
primal high-socle level with \(b\le a\); the dual vectors at their
tops give the same contribution on the dual level. Thus both colors
have dimension \(q_r(b)\) on every primal and dual high-socle level,
inside a level of total dimension \(2q_r(b)\).
Conjugating the coloring by \(g\in C(\widetilde t)\) preserves all
these properties.

We choose \(g\) to ensure the four leaf bounds.
The subspace
\[
 A=\operatorname{Ann}(\im\widetilde u)
\]
lies in \(\ker\widetilde t^*\), because
\(\im\widetilde t\subseteq\im\widetilde u\).
Its dual filtration profile is \(2K_r^T(b)\).
Each dual color has profile \(q_r(b)\).
By Lemmas~\ref{lem:filtered-intersection}
and~\ref{tail:centralizer}, a generic coloring has, in either
dual color, intersection with \(A\) of dimension
\[
 \max\left(0,\max_b\{2K_r^T(b)-q_r(b)\}\right)
 \le 2L_r ,
\]
where the last inequality is Lemma~\ref{tail:threshold-bounds}.
This is exactly the new leaf top: a functional on a color
annihilates the projected incoming image if and only if its extension
by zero to the other color annihilates \(\im\widetilde u\).

Similarly, \(\ker\widetilde v\subseteq\ker\widetilde t\) has profile
\(2K_r^S(b)\).  A generic primal color meets it in dimension at most
\(2M_r\), giving the new leaf socle bound.
Each of these four requirements holds on a nonempty open subset of
\(C(\widetilde t)\).  Since this group is irreducible, all four hold
simultaneously.  Surjectivity onto each of the primal and dual
filtration stabilizers suffices here; the two actions need not be
independently prescribable.

Let \(P_\pm\) be the resulting color projections, and define the new
leaf arrows by
\[
 u_\pm=P_\pm\widetilde u,\qquad
 v_\pm=\widetilde v|_{\widetilde W_\pm}.
\]
Their return maps are zero because
\(P_\pm\widetilde u\widetilde v|_{\widetilde W_\pm}=0\).
At the preceding trunk vertex,
\[
 v_+u_++v_-u_-=
 \widetilde v(P_++P_-)\widetilde u
 =\widetilde v\widetilde u .
\]
The relation, incoming-image sum, and outgoing kernel at that vertex
are therefore unchanged from the doubled chain.  All other trunk
bounds simply double.

It remains to check the dimension vector.  For
\(i<r\), the new trunk dimension is
\(2\sum_{j\le i}\Delta_j\).
The two leaf coefficients for the \(D_{r+1}\)-weight
\[
 (2\Delta_1,\ldots,2\Delta_r,0)
\]
are both \(\sum_{j\le r}\Delta_j=\dim V_r\), precisely the two new
color dimensions.  The new leaf top bounds are \(2L_r\), and the
socle bounds are \(2M_r\), as required.  Applying
Proposition~\ref{prop:preprojective} gives the asserted occurrence.

When \(r=1\), the retained edge is the virtual edge from \(V_0=0\).
Then \(t_1=u_1=v_1=0\), and the construction produces the two
disconnected vertex spaces of \(D_2=A_1\times A_1\).  The same bounds
and dimension calculation apply.
\end{proof}

\subsection{Completion of the saturation theorem}

Figure~\ref{fig:zero-tail-scale} records the changes of rank and
scale in the common-zero-tail branch of the proof. In particular,
the smaller-rank shape is obtained by restricting the original
fundamental columns; it is not required to be a smaller-rank
fundamental-column shape.
\begin{figure}[tb]
\centering
\begin{tikzpicture}[
  box/.style={draw,rounded corners=2pt,align=center,
    text width=3.45cm,minimum height=1.5cm,inner sep=5pt},
  arr/.style={-{Stealth[length=2mm]},line width=.65pt},
  every node/.style={font=\small}]
  \node[box] (original) at (0,0)
    {$D_n$\\tensor occurrence\\
     $(N\lambda,N\mu,N\xi)$};
  \node[box] (doubled) at (4.5,0)
    {$D_{r+1}$\\tensor occurrence\\
     $(2N\widehat\lambda,2N\widehat\mu,2N\widehat\xi)$};
  \node[box] (moment) at (9,0)
    {$D_{r+1}$\\real moments\\
     $(\widehat\lambda,\widehat\mu,\widehat\xi)$};
  \node[box] (smallpath) at (9,-3.2)
    {$D_{r+1}$\\integral schedule\\
     shape $\widehat\eta$, scale $1$};
  \node[box] (largepath) at (4.5,-3.2)
    {$D_n$\\integral schedule\\
     original shape $\eta$, scale $1$};
  \node[box] (conclusion) at (0,-3.2)
    {$D_n$\\tensor occurrence\\
     $(\lambda,\mu,\xi)$};
  \draw[arr] (original)--node[above=30pt,align=center]
    {zero-tail\\removal} (doubled);
  \draw[arr] (doubled)--node[above=30pt,align=center]
    {moments, then\\divide by $2N$} (moment);
  \draw[arr] (moment)--node[left,align=right]
    {integral-schedule\\theorem} (smallpath);
  \draw[arr] (smallpath)--node[below=30pt]
    {pad with zeros} (largepath);
  \draw[arr] (largepath)--node[below=30pt,align=center]
    {column\\intertwiners} (conclusion);
\end{tikzpicture}
\caption{The common-zero-tail branch, with $1\le r\le n-2$.
Here $\eta$ is a fundamental-column shape for the original
$D_n$ moving weight, and $\widehat\eta$ is its restriction to
the first $r+1$ coordinates. Doubling occurs only before passage
to real moments. Both integral schedules and the final tensor
occurrence use the original scale.}
\label{fig:zero-tail-scale}
\end{figure}


\begin{proof}[Proof of Theorem~\ref{thm:main}]
Suppose first that the invariant at scale \(N\in\Z_{>0}\) is nonzero.
Set \(\xi=\nu^*\), so \(V(N\xi)\) occurs in
\(V(N\lambda)\otimes V(N\mu)\), and
\(\lambda+\mu-\xi\) lies in the root lattice.
We may permute the three invariant labels before designating the
moving label \(\mu\).

If one of them has at most one zero coordinate, choose it as \(\mu\).
Lemma~\ref{lem:invariant-moment}, followed by division by \(N\),
gives a real moment solution for \(\lambda,\mu,\xi\).
Choose the shape \(\eta\) made of successive fundamental columns
summing to \(\mu\).
Theorem~\ref{thm:integral-schedule} gives an integral schedule for
this shape, and Proposition~\ref{prop:columns} gives
\(V(\lambda)\otimes V(\mu)\longrightarrow V(\xi)\) nontrivially.
This is the desired invariant.

Otherwise all three labels have at least two zero coordinates.
They are then integer-coordinate weights, nonnegative and supported
on initial coordinate segments.  Choose a label of largest support
as \(\mu\), and denote its support length by \(r\le n-2\).
If \(r=0\), all three representations are trivial.  Assume \(r\ge1\).
Apply Proposition~\ref{prop:zero-tail} to the occurrence at scale
\(N\).  It gives the occurrence in \(D_{r+1}\) for the truncated
labels at scale \(2N\).  Lemma~\ref{lem:invariant-moment}, divided
by \(2N\), therefore gives a real moment solution for the
\emph{unscaled} truncated labels.

Those labels remain integral, and their difference remains in
\(Q(D_{r+1})\): deleting zero coordinates preserves the even
coordinate sum.  The moving truncated label is positive in its first
\(r\) coordinates and has exactly one zero.  Take the fundamental
column shape for the original \(D_n\)-weight \(\mu\).
Since \(M_{r+1}=\cdots=M_n=0\), its fundamental expansion uses
only columns \(e_1+\cdots+e_k\) with \(k\le r\); every increment
therefore already vanishes beyond the retained coordinates.
Restrict each increment to the first \(r+1\) coordinates.
These increments are dominant for \(D_{r+1}\).  They need not be its
fundamental weights; Theorem~\ref{thm:integral-schedule} applies to
this arbitrary dominant-increment shape and gives an unscaled
integral schedule.

Pad this schedule with zero coordinates.  Its roots and even signed
permutations belong to \(D_n\), its endpoint lies in the required
full Weyl orbit, and every switch coefficient and sign is unchanged.
It is therefore an integral schedule for the \emph{original}
fundamental column shape.  Proposition~\ref{prop:columns} again
gives the unscaled invariant.  The use of \(2N\) supplied only real
existence in smaller rank and has introduced no saturation factor.

Conversely, let an unscaled invariant functional be nonzero.
The compact translates of highest-weight vectors span each
irreducible representation, so the functional is nonzero on some
triple of flag vectors.  Its \(N\)-th tensor power is nonzero on the
\(N\)-th powers of those same flag vectors.  Those powers lie in the
Cartan summands of highest weights \(N\lambda,N\mu,N\nu\).
Restriction gives the invariant at scale \(N\).

The arguments include \(D_2=A_1\times A_1\) and \(D_3=A_3\) with
the stated conventions.  If \(\Spin(2)\) is included, it is a torus;
its representations are characters, and the assertion is the
equivalence of the weight-sum equation with its positive integral
multiple.  This completes the proof.
\end{proof}

\begin{remark}[The stretching parameter]\label{tail:fractional}
Positive integral stretching is essential to the formulation.
If arbitrary positive rational \(N\) are allowed whenever all the
scaled weights exist, the unscaled root-lattice condition alone does
not suffice.  For \(D_n\), \(n\ge4\), take
\(\lambda=\mu=\nu=2e_1\).
The triangle contraction
\[
 (x,y)(y,z)(z,x)
\]
of the standard orthogonal form defines a nonzero invariant on
\(V(2e_1)^{\otimes3}\).
Indeed isotropic squares lie in the highest-weight module
\(V(2e_1)\), and, in a hyperbolic basis \(a_i,b_i\) with
\((a_i,b_j)=\delta_{ij}\), the isotropic vectors
\[
 a_1,\qquad b_1,\qquad a_1+b_1+a_2-b_2
\]
have nonzero pairwise products.  At \(N=1/2\) the three labels are
\(e_1\), whose sum \(3e_1\) is outside \(Q(D_n)\); hence the triple
has no invariant.  This does not affect
Theorem~\ref{thm:main}, where \(N\) is a positive integer.
\end{remark}
